% /root/monodromy, 2026-11-19. Root and reduction-descent complete audits passed.
\subsection{Two exact tangent triples in the norm family}
\label{subsec:ss-global-tangent-triples}
\status{Proved and checked. The parent and reduction-descent agents independently
audited the full proof and exact verification programs. The reduction-descent
agent reran both complete global certificates and the separate local verifier.
These are two geometric witnesses, not a new full-image theorem.}

\begin{lemma}[Global witnesses for both affine triple types]
\label{lem:ss-global-tangent-triples}
Use the field, isogeny, monomial tensor basis and kernel elements $G,H$
of Lemma~\ref{lem:ss-global-contact03} and
Lemma~\ref{lem:ss-level3-collision-orbits}.
The following two coefficient vectors define geometrically smooth
source curves:
\[
 \begin{split}
 s_{\mathrm{line}}&=(24,2,49,5,40,45,54,16,46),\\
 s_{\mathrm{nonline}}&=(37,55,55,1,1,2,61,1,5).
 \end{split}
\]
Both are sharp: their Cartier kernels consist exactly of the two
ambient invariant forms. Their norm curves each have one three-branch
singularity consisting of a simply tangent pair and a transverse third
branch, and otherwise only ordinary nodes.
For the first curve the three source points are
\[
 \begin{split}
 P&=((24,23),(56,26)),\\
 Q=P+G&=((51,23),(23,34)),\\
 R=P-G&=((57,25),(55,21)).
 \end{split}
\]
For the second they are
\[
 \begin{split}
 P&=((3,44),(14,44)),\\
 Q=P+G&=((18,34),(50,30)),\\
 R=P+H&=((12,26),(25,53)).
 \end{split}
\]
In both cases $P,Q$ are the tangent pair. In particular these are
witnesses for the affine-line and noncollinear triple orbits,
respectively, and the degree-three relative Hasse divisor is smooth
at both norm curves.
\end{lemma}
\begin{proof}
All coefficients use the integer encoding of
$\mathbf F_2[z]/(z^6+z+1)$ in
Lemma~\ref{lem:ss-global-contact03}. The following exact values
are retained in the discovery and verification outputs:
\[
\begin{array}{c|ccccc}
 &D(s)&Q(s)&c\text{ in }ds(Q)=c\,ds(P)
 &\det(ds(P),ds(R))&\text{contact coefficient}\\\hline
 \mathrm{line}&55&32&31&41&3\\
 \mathrm{nonline}&62&1&35&49&39
\end{array}
\]
In particular $D(s)^2Q(s)\ne0$, so the relative source Cartier map
is bijective. The ambient invariant forms are killed by Cartier and
map isomorphically under the adjunction boundary. Thus the full source
Cartier kernel consists exactly of these two forms.

The local checks are independent of the discovery formula.
The program \path{agents/monodromy/code/verify_tangent_triple_local.py}
reconstructs the finite field by binary polynomial reduction,
checks the translated points and all displayed values, and computes
actual order-two elliptic-curve power series. It chooses a branch
$\gamma(t)$ of the source curve through $P$ to order two and verifies
\[
 s(\gamma(t))=O(t^3),\qquad
 s(\gamma(t)+G)=a t^2+O(t^3),
 \qquad a=3\ \text{or }39.
\]
The nonzero determinant in the table proves that the third branch
is transverse. Hence the intended singularity is the precise
one in Lemma~\ref{lem:ss-tangent-triple-bridge}.

For global smoothness, the same four projective product charts used
in Lemma~\ref{lem:ss-global-contact03} give the unit ideal for the
source singular-point equations, separately for both coefficient
vectors. The norm-collision verifier uses the same eight globally
certified translation matrices. It checks all eight pair-tangency
ideals and all twenty-eight triple-identification ideals on each of
the four charts. Pair-tangency support is exactly $P$ for $G$ and
$Q$ for $-G$, and is otherwise empty.

For the triple ideals, the allowed support is determined before the
Groebner computation: for given $g,h$, retain exactly those
$T\in\{P,Q,R\}$ such that $T+g,T+h\in\{P,Q,R\}$ and that lie
in the chosen chart. If this list is empty, the ideal must be the
unit ideal. For a nonempty list $T_1,\ldots,T_m$, $m\le3$, choose
a linear form $\ell$ separating them. Write $\alpha_i=\ell(T_i)$
and let $p_j$ be the degree-at-most-$m-1$ interpolation polynomial
with $p_j(\alpha_i)=(T_i)_j$. The program reduces
\[
  \prod_i(\ell-\alpha_i)^{32},\qquad
          (x_j-p_j(\ell))^{32}\quad(1\le j\le4)
\]
to zero in the computed Groebner basis. These identities prove
support containment over the algebraic closure: at a common zero,
$\ell=\alpha_i$ for some $i$, and injectivity of the thirty-second
power map then forces every coordinate to equal that of $T_i$.
The prescribed points themselves satisfy the equations. This is
an exact finite-support test, not a search for rational points.

All $144$ collision checks pass for each curve. Thus no additional
triple, quadruple or nontransverse double identification is possible.
The source curves are smooth connected ample divisors and hence
integral. The finite pair-intersection checks exclude translation
invariance, so each source is the normalization of its norm image.
Etaleness of the isogeny excludes any further branch singularity.
This proves both global assertions. Finally
Lemma~\ref{lem:ss-tangent-triple-bridge} applies and proves Hasse
smoothness at both norm curves.
\end{proof}

\begin{remark}[Exact execution record and independent audits]
The discovery program is
\path{agents/monodromy/code/ss_tangent_triple_search.py}.
It uses seeds $202611191140$ and $202611191141$, and accepts the
second and fourth sections at the first base point, respectively.
Its rational orbit and tangency screens only reject candidates.
The complete discovery records are
\path{agents/monodromy/outputs/ss_tangent_triple_collinear.json} and
\path{agents/monodromy/outputs/ss_tangent_triple_noncollinear.json}.
The discovery runtimes were $0.8046595620035077$ and
$0.7408189790003235$ seconds. The first noncollinear attempt had an
extra point in the intended fibre, and the next two had other rational
triple fibres; these rejections are retained in that output.

The global collision program is
\path{agents/monodromy/code/verify_ss_tangent_triple_collisions.py}.
Its full outputs are
\path{agents/monodromy/outputs/verify_ss_tangent_triple_collinear_collisions.json}
and
\path{agents/monodromy/outputs/verify_ss_tangent_triple_noncollinear_collisions.json}.
They passed in $72.36891810200177$ and $69.24018129799515$ seconds.
The separate four-chart source smoothness checks, using the retained
root verifier, passed in $4.423437987999932$ and $3.9984959799985518$
seconds; their complete outputs have the corresponding
\texttt{geometry.json} suffixes in the same directory.
The independent local-series verifier passed in
$0.005125144998601172$ seconds, with output
\path{agents/monodromy/outputs/verify_tangent_triple_local.json}.
The parent independently read the complete proof and passed it. The
reduction-descent agent independently read all programs and the proof,
then reran both $144$-check collision certificates in
$76.3690687260023$ and $74.86146398200071$ seconds. Its four-chart
source checks passed in $3.7411868540002615$ and
$3.3895406299998285$ seconds, and its separate local-series rerun
passed in $0.007957189001899678$ seconds, again producing $3$ and $39$.
The independent outputs are retained in the reduction-descent agent's
output directory with stems
\path{independent_ss_tangent_triple_collinear},
\path{independent_ss_tangent_triple_noncollinear}, and
\path{independent_tangent_triple_local}. The author outputs were
preserved. Both complete independent audits passed.
\end{remark}
